\documentclass[10pt,a4paper]{article}
\usepackage{amsmath,amssymb,amsthm,graphicx}
\usepackage{amsmath,amssymb,amsthm,mathtools,xspace}
\usepackage[no-math]{fontspec}
\usepackage[colorlinks=true,hypertexnames=false]{hyperref}
\newtheorem{theorem}{Theorem}
\newtheorem{lemma}[theorem]{Lemma}
\newtheorem{claim}[theorem]{Claim}
\newtheorem{corollary}[theorem]{Corollary}
\newtheorem{fact}[theorem]{Fact}
\theoremstyle{definition}
\newtheorem{definition}[theorem]{Definition}
\newtheorem{remark}[theorem]{Remark}
\newcommand{\F}{\mathbb C}
\newcommand{\N}{\mathbb N}
\renewcommand{\vec}[1]{\mathbf{#1}}
\newcommand{\va}{\mathbf a}
\newcommand{\vb}{\mathbf b}
\newcommand{\vu}{\mathbf u}
\newcommand{\vx}{\mathbf x}
\newcommand{\vy}{\mathbf y}
\newcommand{\vz}{\mathbf z}
\newcommand{\set}[1]{\left\{#1\right\}}
\newcommand{\abs}[1]{\left|#1\right|}
\DeclareMathOperator{\rank}{rank}
\DeclareMathOperator{\poly}{poly}
\allowdisplaybreaks
\AtBeginDocument{\author{\parbox{.95\textwidth}{\centering Prasad Chaugule, Mrinal Kumar, Nutan Limaye,\\ Chandra Kanta Mohapatra, Adrian She, Srikanth Srinivasan}}}
\title{Schur-family formula hardness conditional on determinant formula hardness}
\author{Prasad Chaugule, Mrinal Kumar, Nutan Limaye, Chandra Kanta Mohapatra, Adrian She, Srikanth Srinivasan}
\date{doi:10.4230/LIPIcs.CCC.2020.14; CC BY 3.0 Unported}
\begin{document}
\maketitle
\noindent\textit{Editorial scope.} One conditional Schur-family formula-hardness/determinant mechanism. Complete composition/Jacobian/symmetric-function support and direct generalized-Vandermonde/skew consequences grouped once;\$\textbackslash{}mathrm\{VF\}\textbackslash{}ne\textbackslash{}mathrm\{VBP\}\$ remains an assumption,not a proved conclusion. No author proof repair or independent certification.\par\medskip
\noindent\textit{Allowed standard prerequisites (editorial).}\begin{itemize}
\item Characteristic-zero algebraic complexity,symmetric-function identities,Jacobian criterion,Taylor expansion and division elimination. Formal field is C;paper remarks extension to characteristiczero or sufficiently large characteristic without a numerical threshold. Standard Jacobian/Taylor/Jacobi–Trudi/division results retained as cited prerequisites;all local author proofs/adaptations retained. Lemma32 is the author's printed proof sketch,including their explicit statement that remaining analogous details are skipped. Full41 formal bibliography entries retained.
\end{itemize}\medskip
\section*{Main statement and conditional consequences}
\setcounter{theorem}{2}
\begin{theorem}[Main Theorem]
\label{thm:main-intro}
Assume that $\lambda= (\lambda_1,\ldots,\lambda_\ell)$ is such that $\lambda_i \geq \lambda_{i+1}+(\ell-1)$ for all $i \in [\ell-1]$, also $\lambda_{\ell} \geq \ell$ and let $d = \sum_i \lambda_i$. Then, for $n \geq \lambda_1 + \ell$, if $s_\lambda(x_1,\ldots,x_n)$ has an algebraic formula of size $s$ and depth $\Delta$, then the $\ell\times \ell$ determinant $(\mathrm{det}_\ell )$ has an algebraic formula of size $\poly(s)$ and depth $\Delta + O(1).$
\end{theorem}


For suitable choices of $\ell,d,n$ above, we can ensure that these parameters are all polynomially related. The theorem then implies that the Schur polynomials do not have algebraic formulas of polynomial size unless the entire complexity class $\mathrm{VBP}$ collapses to the complexity class $\mathrm{VF}$ which consists of polynomials with small formulas. Moreover, the Schur polynomials do not have  \emph{constant-depth} formulas of subexponential size unless the determinant does, a result that would greatly improve the state-of-the-art in this direction~\cite{r16}. 

The above theorem and its proof have several interesting aspects that we now elaborate on. 

\begin{lemma}[Main Technical Lemma (informal)]\label{lem:key-intro}
Let $g_1,\ldots,g_n\in \mathbb{C}[x_1,\ldots,x_n]$ be ``well-behaved'' algebraically independent polynomials. Then, for any homogeneous polynomial $\tilde{f}$, if $f = \tilde{f}(g_1,\ldots,g_n)$ has a formula of size $s$ and depth $\Delta$, the polynomial $\tilde{f}$ has a formula of size $\poly(s)$ and depth $\Delta + O(1).$
\end{lemma}
\paragraph*{Generalized Vandermonde determinants.} The Vandermonde matrix $(x_i^{n-j})_{i,j\in [n]}$ and its determinant are ubiquitous in computation because of their relation to polynomial interpolation. More precisely, the problem of finding a degree-$(n-1)$ univariate polynomial that takes prescribed values at a specified set of $n$ distinct points involves solving a linear system of equations where the underlying matrix is precisely the Vandermonde matrix. It is, therefore, an important fact that the Vandermonde determinant is computationally much easier than the general determinant: in fact, it has the following standard depth-$2$ formula
\[
\mathrm{det}_n\left((x_i^{n-j})_{i,j\in [n]}\right) = \prod_{i,j\in [n]: i < j}(x_i - x_j).
\]
However, it is unclear whether such small expressions continue to exist if we allow the exponents of the variables to vary more generally. For integers $\mu_1 > \mu_2 > \cdots > \mu_n \geq 0$, consider the \emph{generalized Vandermonde matrix} $(x_i^{\mu_j})_{i,j\in [n]}$. Similar to the Vandermonde matrix, the determinant of this matrix is related to the problem of \emph{sparse polynomial interpolation}, where we are trying to interpolate a polynomial only involving  the monomials of degree $\mu_1,\ldots,\mu_n$ through the given points. 

Can we expect that computing any generalized Vandermonde determinant is much easier than computing the determinant itself? It follows from Theorem~\ref{thm:main-intro} and the bialternant formula from Definition~\ref{def: Schur defn} above that the answer to this question is negative: for certain (polynomially large) exponents, the generalized Vandermonde determinant is not much easier than the determinant.




\setcounter{section}{1}

\section{Notations and Preliminaries}
\label{sec:prelims}
Throughout this paper, we assume that we are working over the field $\F$. It is not very hard to see that the results we present can be made to work for fields of characteristic zero or fields of sufficiently large characteristic. Boldface letters are used for tuples of variables e.g. $\vx$ for $(x_1, x_2, \ldots, x_n)$. For $\vb = (b_1, b_2, \ldots, b_n) \in \N^n$ and $\vx = (x_1, x_2, \ldots, x_n)$, we use  $\vx^{\vb}$ to denote  $\prod_{i = 1}^n x_i^{b_i}$. 
We use $|\vb|_1$ to denote $\sum_{i \in [n]} b_i$. %

\subsection{Models of computation}
We start by defining some of the standard models of algebraic computation that we work with in the rest of the paper. 
\begin{definition}[Algebraic circuit]	
	An algebraic circuit $C$ is a directed acyclic graph usually having two kinds of node $+,\times$ with a unique sink vertex having out-degree $0$ called the root. The source vertices having in-degree $0$ are labelled by either field constants or formal variables. Further edges entering into $+$ gate can have field constants on them allowing the plus gate to compute the $\mathbb{F}$-linear combination of its children. The root outputs the polynomial computed by the algebraic circuit. 
\end{definition}

\begin{definition}[Algebraic formula]
	If the underlying graph is a tree instead of a directed acyclic graph then the circuit is called a formula. Otherwise we can define formula to be a circuit having output fan-in at most $1$ for every node.
\end{definition}

\begin{definition}[Algebraic branching program]
	Algebraic branching program (ABP) is a layered graph having a unique source vertex (we call it $s$) and a unique sink vertex (we call it $t$). All the edges are from layer $i$ to $i+1$, and each edge is labelled by a linear polynomial. The weight of a path $p$ is the product of the labels over the edges in $p$. The polynomial that the ABP computes is the sum of all weighted paths from $s$ to $t$. 
\end{definition}

\subsection{Interpolation and Division elimination}
We now state two fairly standard facts about algebraic formula.  The first of these relates the formula complexity of a polynomial to the formula complexity of each of its homogeneous components.  
\begin{lemma}\label{lem:homogeneous components}
Let $P(\vx) \in \F[\vx]$ be a polynomial which can be computed by a formula of size at most $s$ and depth $\Delta$. Then, for every $d$, the homogeneous component of $P$ of degree $d$ can be computed by a formula of size at most $O(s^2)$ and depth $\Delta+O(1)$.  
\end{lemma}
The proof of this lemma is via a standard interpolation argument, where we consider the polynomial $Q(t) = P(x_1t, x_2t, \ldots, x_nt) \in \F(\vx)[t]$ as a univariate in $t$. The point to note is that the homogeneous components of $P$ are coefficients of various powers of $t$ in this new polynomial, and hence can be computed as a linear combination of sufficiently many evaluations of $Q(t)$ for distinct values of $t$ in the base field (which we assume to be large enough). For every $\alpha \in \F$, the formula size of $Q(\alpha)$ is upper bounded by the formula size of $P$. Similarly the depth of the formula of $Q(\alpha)$ is bounded by the depth of $P$. The number of such distinct evaluations needed is upper bounded by one more than the degree of $Q$, which is one more than the degree of $P$ itself. The final observation needed for proving the size upper bound is that a polynomial which can be computed by a formula of size $s$ has degree upper bounded by $s$. Thus, we need to take an appropriately weighted linear combination of $s+1$ distinct substitutions of $t$ in $Q$, each of which has a formula of size at most $s$; thereby giving us an upper bound of $O(s^2)$. Taking linear combinations of such substitutions can be done in depth $O(1)$, which gives the overall depth bound of $\Delta+O(1)$. 

The next statement we need is about the formula complexity of a polynomial which can be written as quotient of two polynomials with small formulas. 
\begin{lemma}\label{lem:division of formulas}
Let $P$ and $R$ be polynomials in $\F[\vx]$ of formula (ABP/circuit) size at most $s$ and depth at most $\Delta$ such that $R$ divides $P$. Then, the polynomial $Q = \frac{P}{R}$ can be computed by a formula (an ABP/circuit resp.) of size at most $\poly(s)$ and depth at most $\Delta+O(1)$. 
\end{lemma}
The proof of this lemma goes via the standard division elimination argument of Strassen and that of~Lemma~\ref{lem:homogeneous components}. We refer the reader to the excellent survey of Shpilka and Yehudayoff~\cite{r37} for formal details on division elimination.

%
%

%
%

\subsection{Algebraic independence and the Jacobian}
The notion of algebraic independence that we now define plays a crucial role in the proofs in the paper. We start with a formal definition. 
\begin{definition}\label{def:algebraic independence}
Polynomials $q_1, q_2, \ldots, q_k \in  \F[\vx]$ are said to be algebraically independent over $\F$ if there is no non-zero polynomial $g(y_1, y_2, \ldots, y_k) \in \F[\vy]$ such that $g(q_1, q_2, \ldots, q_k)$ is identically zero. 
\end{definition}
This definition generalizes the notion of \emph{linear} independence, which is the special case when there is no non-zero polynomial $g$ in $k$ variables and degree $1$ such that $g(q_1, q_2, \ldots, q_k)$ is identically zero. As we shall see next, over fields of characteristic zero ( or sufficiently large characteristic), the notion of algebraic independence is characterized by the rank of the Jacobian matrix defined below. 
\begin{definition}\label{def:jacobian}
The Jacobian matrix of a tuple $(q_1, q_2, \ldots, q_k)$ of $n$ variate polynomials in $\F[\vx]$, denoted by ${\cal J}(q_1, q_2, \ldots, q_k)$ is a $k \times n$ matrix with entries from $\F[\vx]$ whose $(i, j)^{th}$ entry equals $\frac{\partial q_i}{\partial x_j}$. 
\end{definition}
Thus, the row corresponding to $q_i$ in ${\cal J}(q_1, q_2, \ldots, q_k)$ contains all of the $n$ first order partial derivatives of $q_i$. In other words, the $i^{th}$ row of the Jacobian gives us the gradient of $q_i$. The connection between algebraic independence and the Jacobian stems from the following (almost folklore) theorem. 
\begin{theorem}[Jacobian and Algebraic Independence]\label{thm:jac and alg ind}
Let $\left(q_1, q_2, \ldots, q_k \right)$ be a $k$ tuple of $n$ variate polynomials in $\F[\vx]$ of degree at most $d$. Then, $q_1, q_2, \ldots, q_k$ are algebraically independent over $\F$ if and only if the the rank of the Jacobian matrix ${\cal J}(q_1, q_2, \ldots, q_k)$ over the field $\F(\vx)$ is equal to $k$.
\end{theorem}
A proof of this theorem can be found in the survey of Chen, Kayal and Wigderson~\cite[Chapter 3]{r4}. 


\subsection{Taylor's expansion}
\label{sec:prelim-taylor}
For our proof, we need the following well-known form of Taylor's expansion. 
\begin{theorem}\label{thm:taylor}
Let $P \in \F[\vx]$ be an $n$ variate polynomial of degree at most $d$, and  let $\va \in \F^n$ be a point. Then, for an $n$-tuple of variables $\vz$
\[
P(\va + \vz) = \sum_{i = 0}^d \left(\sum_{\vu \in \N^{n}, |\vu|_1 = i} \frac{\vz^{\vu}}{\vu !} \cdot \frac{\partial P}{\partial \vx^{\vu}}\left(\va\right) \right) \, 
\]  
where, for $\vu = (u_1, u_2, \ldots, u_n)$, $\vu! = u_1!\cdot u_2!\cdots u_n!.$
\end{theorem} 
Note that for $i = 0$, the summand $\left(\sum_{\vu \in \N^{n}, |\vu|_1 = i} \frac{\vz^{\vu}}{\vu !} \cdot \frac{\partial P}{\partial \vx^{\vu}}\left(\va\right) \right)$ is just equal to $P(\va)$, and for every positive integer $i$ at most $d$, this summand is a homogeneous polynomial of degree equal to $i$ in $\vz$. Of particular utility to us is the following easy corollary of~Theorem~\ref{thm:taylor}. 
\begin{corollary}\label{cor:taylor truncation}
Let $P \in \F[\vx]$ be an $n$ variate polynomial of degree $d \geq 1$, and  let $\va \in \F^n$ be a point. Then, for an $n$-tuple of variables $\vz$
\[
P(\va + \vz) = P(\va) + \sum_{j = 1}^n z_j \cdot \frac{\partial P}{\partial x_j}(\va) \mod \langle \vz \rangle^2\, .
\]  

\end{corollary}

\subsection{Two useful lemmas}
We use the following (well known) lemma in our arguments. While the lemma is essentially folklore, we sketch a proof for completeness. 
\begin{lemma}\label{lem:remainder}
Let $f(\vx) \in \F[\vx]$ and $P(\vx, y) \in \F[\vx, y]$ be  polynomials such that $P\left(\vx, f(\vx)\right)$  is identically zero. Then, there exists a polynomial $Q(\vx, y) \in \F[\vx, y]$ such that $P(\vx, y) = (y-f(\vx))\cdot Q(\vx, y)$. 
\end{lemma}
\begin{proof}[Proof]
Let $d$ be the degree of $P$ in $y$ and let $C_0(\vx), C_1(\vx), \ldots, C_d(\vx)$ be polynomials in $\F[\vx]$ such that 
\[
P(\vx, y) = \sum_{i = 0}^d C_i(\vx)\cdot y^i \, .
\]
Therefore, $P(\vx, f(\vx))$ can be written as 
\[
P(\vx, f(\vx)) = \sum_{i = 0}^d C_i(\vx)\cdot f(\vx)^i \, .
\]
Subtracting the two expressions above, we get 
\[
P(\vx, y) - P(\vx, f(\vx)) =  \sum_{i = 0}^d C_i(\vx)\cdot y^i  - \sum_{i = 0}^d C_i(\vx)\cdot f(\vx)^i \, .
\]
Now, on the right hand side, the term for $i = 0$ cancels out and on further simplification, we get 
\[
P(\vx, y) - P(\vx, f(\vx)) =  \sum_{i = 1}^d C_i(\vx)\cdot \left(y^i - f(\vx)^i\right)  \, .
\]
Note that for every natural number $i \geq 1$, $y^i-f(\vx)^i$ is divisible by  $(y-f(\vx))$. Therefore, every summand on the right hand side has $(y-f(\vx))$ as a factor, and thus there is a polynomial $Q(\vx, y)$ such that 
\[
P(\vx, y) - P(\vx, f(\vx)) =  (y-f(\vx))\cdot Q(\vx, y)  \, .
\]
Moreover, since $P(\vx, f(\vx))$ is identically zero, we have that $
P(\vx, y) =  (y-f(\vx))\cdot Q(\vx, y)  \, , 
$ thereby completing the proof of the lemma. 
\end{proof}
%
In particular, if $|S| \ge d+1$, then there exists some $\va \in S^n$ satisfying $P(\va) \ne 0$. 
This gives us a brute force deterministic algorithm, running in time $(d+1)^n$, to test if an arithmetic circuit computing a polynomial of degree at most $d$ in $n$ variables is identically zero.


\section{Symmetric polynomials}
\label{sec:prelim-sym-poly}
A polynomial is said to be symmetric if it is invariant under a permutation of variables. We now define some of the families of symmetric polynomials that are discussed in this paper and briefly discuss some of their properties. For a more detailed introduction on symmetric polynomials, we refer the reader to the book~\cite{r28}. We start with the definitions.
\begin{definition}[Elementary symmetric polynomials]\label{def:esym}
The elementary symmetric polynomial of degree $k$ on $n$ variables denoted by $e_k(\vx)$ is defined as follows:
\[
e_k(\vx) = \sum_{S\subseteq [n]} \prod_{i \in S} x_i \, .
\]
\end{definition}
The following fact states a property of the elementary symmetric polynomials which will be useful for our proofs in the later sections. 
\begin{fact}\label{fact:esym and roots}
For all $\alpha_1, \alpha_2, \ldots, \alpha_n \in \F$, if $c_1, c_2, \ldots, c_n$ are field elements such that \[
\prod_{i = 1}^n (z-\alpha_i) = z^{n} - c_1\cdot z^{n-1} + c_2\cdot z^{n-2}- \cdots + (-1)^{n} c_n \, ,
\]
then, for every $j \in [n]$, $c_j = e_j(\alpha_1, \alpha_2, \ldots, \alpha_n)$.
\end{fact}

\begin{definition}[Homogeneous (Complete) symmetric polynomials] The homogeneous symmetric polynomial of degree $k$ on $n$ variables denoted by $h_k(\vx)$ is defined as follows:
	$$h_k(\vx)=\sum_{\vb \in \N^n: |\vb|_1 = k} \vx^\vb.$$ 
\end{definition}

\begin{definition}[Homogeneous (Complete) symmetric polynomials] The homogeneous symmetric polynomial of degree $k$ on $n$ variables denoted by $h_k(\vx)$ is defined as follows:
	$$h_k(\vx)=\sum_{\vb \in \N^n: |\vb|_1 = k} \vx^\vb.$$ 
\end{definition}

\begin{definition}[Power symmetric polynomials] The power symmetric polynomial of degree $k$ on $n$ variables denoted by $p_k(\vx)$ is defined as follows:
	$p_k= \sum_{i=1}^n  x_i^k$.
\end{definition}

These sets of polynomials are algebraically independent. The following fact states this formally. 
\begin{fact}\label{fact:sym ad}
Let $\vx$ be an $n$ tuple of variables. Then, elementary symmetric polynomials $e_1(\vx), \ldots, e_n(\vx)$   are algebraically independent over $\F$. Similarly,  homogeneous symmetric polynomials $h_1(\vx), \ldots, h_n(\vx)$ and power symmetric polynomials $p_1(\vx), p_2(\vx), \ldots, p_n(\vx)$ are also algebraically independent over $\F$. 
\end{fact}

%

We now formally state the fundamental theorem of symmetric polynomials, which essentially says that over field of characteristic zero, every symmetric polynomial can be written as a unique polynomial combinations of the elementary symmetric polynomials (similarly, for power symmetric polynomials or homogeneous symmetric polynomials).
\begin{theorem}[The fundamental theorem of symmetric polynomials]
\label{thm:fundamental}
For a symmetric polynomial $f_{sym} \in \F[\vx]$ there exists a unique polynomial $f \in \F[\vx]$ s.t $f_{sym}=f_E(e_1(\vx),e_2(\vx), \ldots e_n(\vx))$ where $e_i(\vx)$ is the elementary symmetric polynomial of degree $i$. 
 
Similarly, there exists a unique polynomial $f_H \in \F[\vx]$ such that $f_{sym} =f_H(h_1(\vx),h_2(\vx), \ldots h_n(\vx))$ and a unique polynomial $f_P \in \F[\vx]$ such that $f_{sym} =f_P(p_1(\vx),p_2(\vx), \ldots p_n(\vx))$, where $h_i(\vx)$ is the homogeneous symmetric polynomial of degree $i$ and $p_i(\vx)$ is the power symmetric polynomial of degree $i$. 
\end{theorem} 

\subsection{Schur polynomials}
\label{sec:prelim-schur}
A partition of a natural number $d$ is any sequence $\lambda = (\lambda_1,\lambda_2 \ldots, \lambda_{\ell} )$ of non-negative integers in a non-increasing order $\lambda_1 \geq \lambda_2 \ldots \geq  \lambda_{\ell} \geq 0$ such that $\sum_{i = 1}^{\ell} \lambda_i = d$.\footnote{Usually the $\lambda_i$s are assumed to be positive integers as defined earlier. For the sake of notational convenience we allow trailing zeroes in the definition of $\lambda$.}
The number of non-zero parts of $\lambda$ is called the length of $\lambda$ and is denoted by $l(\lambda)$. The weight of $\lambda$, denoted by $\abs{\lambda}$ is defined to be the sum of each individual component, i.e. $|\lambda|=\lambda_1 +
\lambda_2+ \cdots + \lambda_{l(\lambda)}$. If $|\lambda|=d$, then we say that $\lambda$ is a partition of the number $d$ or alternatively a partition of \emph{degree} $d$.  


Let $\lambda$ be a partition of the number $d$. A Ferrers diagram (or simply a diagram) of shape $\lambda$ is is a left-aligned two-dimensional array of boxes with the $i$th row containing $\lambda_i$ many boxes.  The conjugate of $\lambda$, denoted by $\lambda'$, is the diagram obtained by switching the rows and columns of the diagram of $\lambda$. 

%
%
 
\begin{definition}[Schur polynomials] \label{def: Schur defn}
Let $\lambda$ be a partition of degree $d$ and let $l(\lambda) \leq n$. Then the Schur polynomial $s_\lambda(\vx)$ is defined as 
	
	$$s_\lambda(\vx)= \frac{a_{\lambda+\delta}(\vx)}{a_\delta(\vx)}$$
	where $$\delta = (n-1,n-2, \ldots \ldots 2,1,0)$$ \quad  $$a_{\lambda+\delta}(\vx)=\det(x_i^{\lambda_j+n-j})_{1 \leq i,j \leq n}$$ \quad 
	$$a_{\delta}(\vx)=\det(x_i^{n-j})_{1 \leq i,j \leq n}= \prod\limits_{\substack{1 \leq i<j \leq n}}(x_i-x_j)$$
\end{definition}

We first observe that $s_\lambda(\vx)$ is a symmetric polynomial. To see this, note that if $x_i=x_j$ for any $i \neq j$ then $a_{\lambda+\delta}(\vx)$ is $0$. Thus, by Lemma~\ref{lem:remainder} and the fact that $x_i-x_j$ and $x_{i'}-x_{j'}$ do not share a common factor unless $\set{i, j} = \set{i', j'}$, $\prod_{i<j}(x_i-x_j)$ is factor of $a_{\lambda+\delta}(\vx)$ i.e., $a_\delta(\vx)$ is a factor of $a_{\lambda+\delta}(\vx)$. Therefore, $s_\lambda(\vx)$ is a polynomial. Moreover, for any permutation of variables, the sign changes in the numerator and the denominator are the same, and thus their ratio does not see a change in sign. This implies that $s_{\lambda}(\vx)$ is a symmetric polynomial. 

We now state the classical Jacobi-Trudi identities which relates Schur polynomials to the elementary symmetric and homogeneous symmetric polynomials.
%

%
\begin{theorem}[Jacobi-Trudi identities] \label {thm:jacobi-trudi}  \leavevmode
	
	\begin{enumerate}
		\item[(1)] $s_\lambda(\vx)=\det(h_{\lambda_i -i+j}(\vx))_{1\leq i,j \leq \ell}$, where $\lambda= (\lambda_1, \ldots, \lambda_\ell)$.
		
		\item[(2)] $s_\lambda(\vx)=$ $\det(e_{\lambda'_i -i+j}(\vx))_{1\leq i,j \leq m}$,  where $\lambda'$ is the conjugate of $\lambda$ and $m= l(\lambda')$ . 
	\end{enumerate} 
	In particular, 
	\vspace{0.3cm}
\begin{align*}
s_\lambda &=  \begin{vmatrix}h_{\lambda_1}&h_{\lambda_1+1}&\cdots&h_{\lambda_1+\ell-1}\\h_{\lambda_2-1}&h_{\lambda_2}&\cdots&h_{\lambda_2+\ell-2}\\\vdots&\vdots&\ddots&\vdots\\h_{\lambda_\ell-\ell+1}&h_{\lambda_\ell-\ell+2}&\cdots&h_{\lambda_\ell}\end{vmatrix}_{\ell \times \ell} \\
	&= \begin{vmatrix}e_{\lambda'_1}&e_{\lambda'_1+1}&\cdots&e_{\lambda'_1+m-1}\\e_{\lambda'_2-1}&e_{\lambda'_2}&\cdots&e_{\lambda'_2+m-2}\\\vdots&\vdots&\ddots&\vdots\\e_{\lambda'_m-m+1}&e_{\lambda'_m-m+2}&\cdots&e_{\lambda'_m}\end{vmatrix}_{m\times m}
\end{align*}
Note that the identity depends only on the properties of $\lambda$ (i.e. $\ell$ or $m$) and does not depend on the number of variables $n$. 
\end{theorem}

%
\begin{theorem}
For any $\lambda$, $s_\lambda(\vx)$ can be computed using a small ABP, hence by a small algebraic circuit.
\end{theorem} 

\begin{proof}
For polynomials $P,Q \in \mathbb{C}[x_1,x_2 \ldots x_n]$ such that both $P$ and $Q$ have small ABPs, then by Lemma~\ref{lem:division of formulas} $R=\frac{P}{Q}$ also has a small ABP. Homogenization or interpolation can be used to extract the required polynomial without much blow up. Here both $a_{\lambda+\delta}(\vx)$, $a_\delta(\vx)$ have small ABPs (as they are small determinants), thus $s_\lambda(\vx)$ has an ABP of polynomial size which also implies that is has an algebraic circuit of polynomial size.
\end{proof}

It is well-known that $a_\delta(\vx)$, also known as the Principal Vandermonde Determinant, has a small algebraic formula. 
However much less is known about the complexity of $a_{\lambda+\delta}(\vx)$. These polynomials are also known as Generalized Vandermonde determinants and are well studied (see for instance~\cite{r18}). %
To the best of our knowledge, before this work it was not known whether for all $\lambda$, $a_{\lambda+\delta}(\vx)$s  have small formulas. Suppose they did, then by Lemma~\ref{lem:division of formulas}, we get that $s_\lambda(\vx)$ also have small formulas for all $\lambda$. 
In this paper we show that there exists some $\lambda$ for which $s_\lambda(\vx)$ does not have a small formula unless the Determinant has a small formula. This in particular implies that there exist $\lambda$ such that  $a_{\lambda+\delta}(\vx)$ cannot be computed using small formulas (unless the Determinant can be computed by a small formula).

%
%
%
%
%

\section{Proofs of main results}
\subsection{Proof of Lemma~\ref{lem:key-intro}}\label{sec:key lemma}
We start with the following definition, which is crucial for our proofs. 
\begin{definition}[Property S]\label{def:prop s}
A set of $n$ variate polynomials $\{q_1, q_2, \ldots, q_k\} \subseteq \F[\vx]$ is said to satisfy Property S, if there exists an $\va \in \F^n$ such that
\begin{itemize}
\item For all $i \in [k]$, $q_i(\va) = 0$, and,
\item The rank of the Jacobian matrix of $q_1, q_2, \ldots, q_k$ when evaluated at $\va$ is equal to its symbolic rank, i.e. ${\rank}_{\F}\left({\cal J}(q_1, q_2, \ldots, q_k)(\va)\right) = \rank_{\F(\vx)}\left({\cal J}(q_1, q_2, \ldots, q_k)\right)$.
\end{itemize} 
\end{definition}
Property S gives us a concrete way to capture an appropriate notion of \emph{niceness} of a set of algebraically independent polynomials. The following lemma which uses this notion is a key technical ingredient of our proofs. 
\begin{lemma}\label{lem:key}
Let $\{q_1, q_2, \ldots, q_k\} \in \F[x_1, x_2, \ldots, x_n]$ be a set of algebraically independent polynomials which satisfy Property S. Let $g \in \F[z_1, z_2, \ldots, z_k]$ be a homogeneous $k$ variate polynomial of degree equal to $d$ such that the composed polynomial $g(q_1, q_2, \ldots, q_k) \in \F[\vx]$ has an algebraic formula of size $s$ and depth $\Delta$. Then, $g(z_1, z_2, \ldots, z_k)$ can be computed by an algebraic formula of size $O(s^2n)$ and depth $\Delta+O(1)$.  
\end{lemma}

\begin{proof}
Let $\Phi$ be the formula of size $s$ which computes the polynomial $g\left(q_1(\vx), q_2(\vx), \ldots, q_k(\vx)\right)$. Since $q_1, q_2, \ldots, q_k$ satisfy Property S, there is an $\va \in \F^n$ such that $q_1(\va) = q_2(\va) = \cdots = q_k(\va) = 0$ and ${\rank}_{\F}\left({\cal J}(q_1, q_2, \ldots, q_k)(\va)\right) = \rank_{\F(\vx)}\left({\cal J}(q_1, q_2, \ldots, q_k)\right)$. Moreover, since they are algebraically independent, the rank of $\left({\cal J}(q_1, q_2, \ldots, q_k)(\va)\right)$ is equal to $k$. Thus, $${\rank}_{\F}\left({\cal J}(q_1, \ldots, q_k)(\va)\right) = k \, .$$
Applying~Corollary~\ref{cor:taylor truncation} to each $q_i(\vx)$ around this point $\va \in \F^n$, we get  
\begin{align*}
q_i(\va + \vx) &= \sum_{j = 1}^n x_j\cdot \frac{\partial q_i}{\partial x_j} (\va)  \mod \langle \vx \rangle^2 \, .
\end{align*}
Observe that $i^{th}$ row of the matrix $\left({\cal J}(q_1, q_2, \ldots, q_k)(\va)\right)$ is the vector $\left(\frac{\partial q_i}{\partial x_1} (\va), \frac{\partial q_i}{\partial x_2} (\va), \cdots, \frac{\partial q_i}{\partial x_n} (\va) \right)$ and by the choice of $\va$, these vectors  are linearly independent. Thus, the homogeneous linear forms $u_1(\vx), u_2(\vx), \ldots, u_k(\vx)$ are linearly independent, where $u_i(\vx)$ is defined as 
\[
u_i(\vx) = \sum_{j = 1}^n x_j\cdot \frac{\partial q_i}{\partial x_j} (\va)\, .
\]
The rest of the proof follows immediately from the following two claims. We state the claims and use them to complete the proof of this lemma, and then  move on to prove the claims. 
\begin{claim}\label{claim:using homogeneity of g}
Let $d$ be the degree of $g(\vx)$. Then, the homogeneous component of degree $d$ of the polynomial $g(q_1, q_2, \ldots, q_k)$ is equal to $g(u_1, u_2, \ldots, u_k)$. 
\end{claim}

\begin{claim}\label{claim:invertible linear transformation}
If $g(u_1, u_2, \ldots, u_k)$ has a formula of size $s'$ and depth $\Delta$, then $g(\vz)$ has a formula of size at most $s'n$ and depth $\Delta+O(1)$. 
\end{claim}
To complete the proof of the lemma, observe that given the formula $\Phi$ of size at most $s$ and depth at most $\Delta$ which computes $g(q_1, q_2, \ldots, q_k)$, we know from~Lemma~\ref{lem:homogeneous components} that the homogeneous component of degree $d$ of  $g(q_1, q_2, \ldots, q_k)$ can be computed by a formula $\Phi_1$ of size at most $O(s^2)$ and depth at most $\Delta+O(1)$. Moreover, from the homogeneity of $g$ and Claim~\ref{claim:using homogeneity of g}, we also know that $\Phi_1$ computes the polynomial $g(u_1, u_2, \ldots, u_k)$, where $u_1, u_2, \ldots, u_k$ are linearly independent linear forms. Thus, from Claim~\ref{claim:invertible linear transformation}, this implies that $g(z_1, z_2, \ldots, z_k)$ has a formula of size at most $O(s^2n)$ and depth $\Delta+O(1)$. This completes the proof of the lemma, modulo the two claims which we prove next. %
\end{proof}

\begin{proof}[Proof of Claim~\ref{claim:using homogeneity of g}]
Let $f_1, f_2, \ldots, f_k \in \F[\vx]$ be polynomials which are zero modulo $\langle \vx \rangle^2$ (i.e. every monomial in $f_1, f_2, \ldots, f_k$ has degree at least $2$) such that for every $i$, $q_i(\va + \vx) = u_i(\vx) + f_i(\vx)$. Since $g$ is a  homogeneous polynomial of degree $d$, it can be expressed as 
\[
g(\vz) = \sum_{\vb \in \N^k, |\vb|_1 = d} \alpha_{\vb} \vz^{\vb} \,  ,
\] for field constants $\alpha_{\vb}$. 
%
  
Let $\vb \in \N^k$ be any vector such that $|\vb|_1 = d$. Observe that the homogeneous component of degree  $d$ of the polynomial 
\[
\prod_{j = 1}^k q_j(\va + \vx)^{b_j} = \prod_{j = 1}^k (u_j(\vx) + f_j(\vx))^{b_j}
\] is equal to $\prod_{j = 1}^k u_j^{b_j} \, .
$
  Thus, by linearity, the homogeneous component of degree $d$ of  \[
  g\left(q_1(\va + \vx), q_2(\va + \vx), \ldots, q_k(\va + \vx)\right) = \sum_{\vb \in \N^k, |\vb|_1 = d} \alpha_{\vb} \cdot \prod_{j = 1}^k (q_j(\va + \vx ))^{\vb_j}\]
   equals \[
   \sum_{\vb \in \N^k, |\vb|_1 = d} \alpha_{\vb}\cdot  \prod_{j = 1}^k (u_j)^{\vb_j} \, , 
   \] which, in turn is the equal to the polynomial $g(u_1, u_2, \ldots, u_k)$.
\end{proof} 
%
\begin{proof}[Proof of Claim~\ref{claim:invertible linear transformation}]
The idea for the proof of this claim is to show that each variable $x_j$ can be replaced by a homogeneous linear form $\ell_j(\vz)$ in the variables $\vz$ such that for every $i \in [k]$, the linear form $u_i$ satisfies $u_i(\ell_1(\vz), \ell_2(\vz), \ldots, \ell_n(\vz)) = z_i$. Thus, under this linear transformation, the composed polynomial $g(u_1(\vx), u_2(\vx), \ldots, u_k(\vx)) \in \F[\vx]$ gets mapped to the polynomial $g(z_1, z_2, \ldots, z_k)$.  Once we can show an existence of these linear forms $\ell_1, \ell_2, \ldots, \ell_n$, the bounds on the formula size and depth follow immediately since all we need to do to obtain a formula for $g(\vz)$ is to replace every occurrence of a variable in $\vx$, e.g $x_j$ by the linear form $\ell_j(\vz)$. Since every such linear form has a formula of size at most $k$ and depth $O(1)$, this process blows up the formula size by a multiplicative factor of at most $O(k)$ and the depth by an additive factor of $O(1)$. 

Intuitively, to obtain these linear forms, we just \emph{solve} the system of linear equations  $U \cdot \vx^{T} = \vz^T$, where $U$ is the $k \times n$ matrix whose $i^{th}$ row is $u_i$.  Since  the rank of $U$ is equal to $k$,  let $U'$ be an invertible $k \times k$ submatrix of $U$, and let $V$ be the inverse of $U'$ and let $v_1, v_2, \ldots, v_k$ be the rows of $V$. Moreover, for brevity, let us assume that $U'$ consists of the first $k$ columns of $U$. Observe that $(V\cdot U)$ is a $k \times n$ matrix such that its leftmost $k \times k$ sub-matrix is the identity matrix. We are now ready to define the linear forms $\{\ell_j : j \in [n]\}$.  For $j \in [k]$, let $\ell_j(\vz)$ be equal to $v_j(\vz)$ and for $j > k$, $\ell_j(\vz)$ is defined to be zero. It is straightforward to check that this definition satisfies the desired property and we skip the details. 
\end{proof}
%
\subsection{Formula complexity of Schur polynomials}
\label{sec:schur-formula-complexity}
%
%
%
%
%
%
%
%
%
%
%
%
%
%
We are now ready to prove Theorem~\ref{thm:main-intro}. Our first stop, which we reach in the next two lemmas, is to show that which shows that the elementary symmetric polynomials of degree at most $n-1$ on $n$ variables are \emph{well behaved}. We start by establishing a sufficient condition for their Jacobian matrix to have full rank at a point. 
\begin{lemma}\label{lem:jacobian of esyms}
Let $\vx = (x_1, x_2, \ldots, x_n)$ be an $n$-tuple of variables. Let $J(\vx)$ be the Jacobian of $e_1(\vx), e_2(\vx), \ldots, e_{n-1}(\vx)$, and let $\vb = (b_1, b_2, \ldots, b_n) \in \F^n$ be such that for all $i \neq j$, $b_i \neq b_j$. Then, 
\[
\rank_{\F(\vx)}\left(J(\vx) \right) = \rank_{\F}(J(\vb)) = n-1 \, .
\]
\end{lemma}
\begin{proof}
Using Fact~\ref{fact:sym ad}, we can observe that  the $n$-variate polynomials $e_1(\vx), e_2(\vx), \ldots, e_{n-1}(\vx)$ are algebraically independent.
Therefore, from Theorem~\ref{thm:jac and alg ind} we know that $\rank_{\F(\vx)}\left(J(\vx) \right)$ is equal to $n-1$. 

Let $J'(\vx)$ be any $n-1 \times n-1$ submatrix of $J(\vx)$ of rank equal to $n-1$ and by symmetry we can assume that the columns in $J'(\vx)$ come from the first $n-1$ columns of $J(\vx)$. Thus, for $i, j \in [n-1]$, the $(i, j)$ entry of $J'(\vx)$ is equal to $\frac{\partial e_i}{\partial x_j}$. We now show that the rank of $J'(\vb)$ over $\F$ is equal to $n-1$ and this would complete the proof. To this end, we observe that the determinant of $J'(\vx)$ is a non-zero scalar multiple of $\prod_{i, i' \in [n-1], i \neq i'}(x_i-x_{i'})$. Since the coordinates of $\vb$ are all distinct, this determinant remains non-zero on $\vb$. 

From the definition of $J'(\vx)$, we know that the entries in its $i^{th}$ row are homogeneous polynomials of degree equal to $i-1$. Thus, $\det(J'(\vx))$ is a homogeneous polynomial in $\vx$ of degree equal to \[
0 + 1 + 2 + \cdots + n-2 = \frac{(n-2)(n-1)}{2} \, . \] Recall that the $(i, j)$ entry of $J'(\vx)$, is equal to $\frac{\partial e_i}{\partial x_j}$, which is equal to $\sum_{S \subseteq [n]/ \{i\}}\prod_{k \in S} x_k$. Thus, if we replace every occurrence of the variable $x_i$ by the variable $x_{i'}$ for $i \neq i' \in [n-1]$, then columns $i, i'$ in $J'(\vx)$ become identical, and hence $\det(J'(\vx))$ is identically zero. Thus, by Lemma~\ref{lem:remainder}, $(x_i-x_{i'})$ is a factor of $\det(J'(\vx))$. Also, for any two distinct sets $\{i_1, i_1'\}$ and $\{i_2, i_2'\}$ where $i_1 \neq i_1'$ and $i_2 \neq i_2'$, the polynomials $(x_{i_1}-x_{i_1'})$ and $(x_{i_2}-x_{i_2'})$ do not share a non-trivial divisor. Thus, the determinant of $J'(\vx)$ must be divisible by the polynomial $\prod_{i, i' \in [n-1], i \neq i'}(x_i-x_{i'})$. Moreover, we observed that the degree of determinant of $J'(\vx)$ is equal to $\frac{(n-2)(n-1)}{2}$, which is also equal to the degree of $\prod_{i, i' \in [n-1], i \neq i'}(x_i-x_{i'})$. Thus, they must be non-zero scalar multiples of each other. This observation, together with the fact that the coordinates of $\vb$ are all distinct, shows that $\det(J'(\vx))$ is non-zero at $\vb$ and hence, $\rank(J(\vb))$ equals $n-1$ over $\F$.
\end{proof}
We now use Lemma~\ref{lem:key} to show that $e_1(\vx), e_2(\vx), \ldots, e_{n-1}(\vx)$ satisfy Property S, where $n$ is the number of variables.
\begin{lemma}\label{lem:esym prop s}
Let $\vx = (x_1, x_2, \ldots, x_n)$ be an $n$-tuple of variables. Then, the set of elementary symmetric polynomials of degree at most $n-1$ on $\vx$, i.e. the set $\{e_1(\vx), e_2(\vx), \ldots, e_{n-1}(\vx)\}$ of polynomials satisfies Property S. 
\end{lemma}

\begin{proof}
Let $\va = (1, \omega, \omega^2, \ldots, \omega^{n-1})$, where $\omega$ is the primitive $n^{th}$ root of unity. So, we have the following identity 
\[
z^n-1 = \prod_{i=1}^{n} (z-\omega^{i-1}) \, .
 \] 
However, from ~Fact~\ref{fact:esym and roots}, we also know that $z^n - 1$, which equals $\prod_{i = 1}^n (z-a_i)$ can be expressed as
\[
\prod_{i = 1}^n (z-a_i) = z^{n} - c_1\cdot z^{n-1} + c_2\cdot z^{n-2}- \cdots + (-1)^{n} c_n \, ,
\]
where $c_i = e_i(\va)$.  
Comparing these two expressions for $z^n -1$, we get that for all $1\leq i < n$, $c_i = e_i(\va)=0$ . Thus, $\va=(1, \omega, \omega^2, \ldots, \omega^{n-1})$ is a common zero of $e_1, e_2, \ldots, e_{n-1}$, which satisfies the first item in Definition~\ref{def:prop s}. Moreover, since $\omega$ is a primitive $n^{th}$ root of one, we also know that for all pairs $i \neq j$, $a_i \neq a_j$. Therefore, by Lemma~\ref{lem:jacobian of esyms}, $\va$ satisfies the second condition in Definition~\ref{def:prop s}. Thus, the $n$ variate polynomials $e_1(\vx), e_2(\vx), \ldots, e_{n-1}(\vx)$ satisfy Property S.  
\end{proof}

%
%
%
%
%
%
%
%
%
%
%
%
%
We now observe that analogous statements are also true for homogeneous symmetric polynomials and power symmetric polynomials as well.  
\begin{lemma}\label{lem:hsym psym prop S} 
Let $\vx = (x_1, x_2, \ldots, x_n)$ be an $n$-tuple variables. Then the set of complete symmetric polynomials of degree at most $n-1$ on $\vx$, i.e. the set $\{h_1(\vx), h_2(\vx), \ldots, h_{n-1}(\vx)\}$ of polynomials satisfies Property S. Similarly, the set $\{p_1(\vx), p_2(\vx), \ldots, p_{n-1}(\vx)\}$ of power symmetric polynomials of degree at most $n-1$ also satisfies Property S.
\end{lemma}

\begin{proof}[Proof sketch]
	The proof goes via  generating functions of $e_i$, $h_i$, $p_i$ and the relations among them. We denote the generating function of $e_i$, $h_i$, $p_i$ by $E(t),H(t),P(t)$ respectively. The following relations are known between these polynomials. (See for instance~\cite{r28}.)
	$$E(t)= \prod_{i \geq 1} (1+x_it)=\sum_{k \geq 0} e_kt^k$$ 
	$$H(t)=\prod_{i \geq 1} \frac{1}{1-x_it}=\sum_{k \geq 0}  h_kt^k$$
	\quad $$ \quad P(t)=\prod_{i\geq 1}\frac{x_i}{1-x_it}= \sum_{k \geq 1}  p_kt^{k-1}$$
It is easy to see that $E(-t)H(t)=1$.  Therefore we get, 
\begin{equation} \label{equ: h to e gen fun}
H(t)= \frac{1}{E(-t)} =\frac{1}{1-e_1t+e_2t^2 \ldots +(-1)^ne_nt^n}
\end{equation}
From Lemma~\ref{lem:esym prop s}, we know there exists a point $\va=(a_1, a_2, \ldots, a_n)$, where $e_1(\va), e_2(\va), \ldots, e_{n-1}(\va)$ vanish and $e_n(\va)$ is non-zero. We evaluate the above equation at the same $\va$. We denote the above evaluation by $H(t)|_{\vx=\va}$.
$$H(t)|_{\vx=\va}=1- (-1)^ne_n(\va)t^n +((-1)^ne_n(\va)t^n)^2 \ldots $$ 
	
Observing the equation for $H(t)|_{\vx=\va}$, it follows that at point $\va$ , $h_1(\va), h_2(\va), \ldots, h_{n-1}(\va)$ are zero and $h_n(\va)$ is non-zero. %
	
An analogous relation between $P(t)$ and $E(t)$, given by
	  \[
	  P(t)=\frac{E^\prime(-t)}{E(-t)} \, ,
	  \] 
where $E^\prime(-t)$ is the first derivative of $E(-t)$ with respect to $t$, can be used to show that the power symmetric polynomials $p_1(\va),\ldots, p_{n-1}(\va)$ are zero and  $p_n(\va)$ is non-zero. 
%
Thus, we are halfway towards showing that the set of power symmetric polynomials of degree at most $n-1$ and homogeneous symmetric polynomials of degree at most $n-1$ also satisfy Property S. It remains to be argued that the rank of the Jacobian of these polynomials at $\va$ is equal to $n-1$. The proof is analogous to the argument for the elementary symmetric polynomials, as in the proof of Lemma~\ref{lem:esym prop s}. We skip the remaining details for brevity. 
\end{proof}
From Definition~\ref{def:prop s} of Property S and Theorem~\ref{thm:jac and alg ind}, it follows that is a set $A$ of polynomials satisfies Property S, then all the non-empty subsets of $A$ also satisfy Property S. Thus, we have the following corollary of Lemma~\ref{lem:esym prop s} and Lemma~\ref{lem:hsym psym prop S}.
\begin{corollary}\label{cor:prop s subsets}
Let $i_1 < \dots < i_k < n$ be positive integers. Consider the sets $E = \{e_{i_1}(\vx), \dots, e_{i_k}(\vx)\}$ and $H = \{h_{i_1}(\vx), \dots, h_{i_k}(\vx)\}$ of elementary and homogenous symmetric polynomials respectively in $n > i_k$ variables. Then both $E$ and $H$ satisfy property $S$. 
\end{corollary}

%
%
%


%
%
%
%

We are now ready to prove the main theorem. 

\begin{theorem}[Main theorem]
\label{thm:main-intro proof}
Let $\lambda= (\lambda_1,\ldots,\lambda_\ell)$ be a partition of $d$ such that  $\lambda_i \geq  \lambda_{i+1}+(\ell-1)$ for all $i \in [\ell-1]$, and $\lambda_{\ell} \geq \ell$. Then, for all $n$, such that $n \geq \lambda_1 + \ell$, if $s_\lambda(x_1,\ldots,x_n)$ has an algebraic formula of size $s$ and depth $\Delta$, then the $\ell\times \ell$ determinant $(\mathrm{det}_\ell )$ has an algebraic formula of size $\poly(s)$ and depth $\Delta + O(1).$  
\end{theorem}

\begin{proof}
	We use Jacobi-Trudi Identity from Theorem~\ref{thm:jacobi-trudi}, which expresses $s_\lambda(\vx)$ in the form of homogeneous symmetric determinant. 
\begin{align*}
s_\lambda &=  \begin{vmatrix}h_{\lambda_1}&h_{\lambda_1+1}&\cdots&h_{\lambda_1+\ell-1}\\h_{\lambda_2-1}&h_{\lambda_2}&\cdots&h_{\lambda_2+\ell-2}\\\vdots&\vdots&\ddots&\vdots\\h_{\lambda_\ell-\ell+1}&h_{\lambda_\ell-\ell+2}&\cdots&h_{\lambda_\ell}\end{vmatrix}
\end{align*}	
Let $M_{\lambda}$ denote the matrix on the right hand side in the above equality. By our choice of $\lambda$, observe that the highest degree entry of $M_{\lambda}$ is $h_{\lambda_1 + \ell -1}$, which has degree at most $n-1$, and the lowest degree entry of $M_{\lambda}$ is $h_{\lambda_{\ell}-\ell + 1}$ which for $\lambda_{\ell} \geq \ell$ has degree at least $1$. Moreover, from the choice of $\lambda$, it also follows that all the entries of $M_{\lambda}$ are distinct.
%
%
 %

From Lemma~\ref{lem:hsym psym prop S}, we know that there exists a point $\va$ for which the $n$-variate polynomials $\{h_1({\vx}),h_2(\vx), \ldots h_{n-1}(\vx)\} $ satisfy the Property $S$, where $n$ can be taken to be strictly greater than $\ell^2$. Thus, now if we take the $\ell^2$-variate homogeneous polynomial $g$ to be the symbolic determinant of an $\ell \times \ell$ matrix, then $s_{\lambda}$ is obtained by a composition of $g$ with a subset of polynomials $h_1, h_2, \ldots, h_{n-1}$. 
%
%
 
Thus, by Lemma~\ref{lem:key}, we get that if $g(h_1({\vx}),h_2(\vx), \ldots h_{\ell^2}(\vx))$ (i.e $s_\lambda$) has an algebraic formula of size $s$ and depth $\Delta$, then $g(x_1,x_2, \ldots x_{\ell^2})$ also has a small formula of size $\poly(s)$ and depth $\Delta+O(1)$ . 
%
\end{proof} 
%
%


%

%
%
%
%
%
%
%

\begin{remark}[Contrasting hard and easy $\lambda$s]
\label{rem:easy-hard-lambda}
Let $\lambda = (\ell^2, \ell^2-\ell, \ldots, 2\ell, \ell, 0, 0, \ldots, 0,0,0)$ and let $\tilde{\lambda} = ( n\ell, (n-1)\ell, \ldots, \ell^2,\ell^2-\ell, \ldots, 3\ell, 2\ell, \ell)$, where we will take $n \geq \ell^2+\ell$ and $\ell>0$. It is easy to see that $\tilde{\lambda}$ forms an arithmetic progression in which the difference between the successive terms is $\ell$. Whereas $\lambda$ is a truncated arithmetic progression, in which  $n-\ell$-many elements are zeroes and the non-zero elements form an arithmetic progression. 

Here the structure of $\lambda$ and $\tilde{\lambda}$ is quite similar, but the algebraic complexities of $s_\lambda$ and $s_{\tilde{\lambda}}$ are  different. 
In particular, Theorem~\ref{thm:main-intro proof} is applicable to $\lambda$. Therefore we can conclude that $s_\lambda$ does not have a small algebraic formula unless the determinant has a small algebraic formula. 

On the other hand, there are small formulas for $s_{\tilde{\lambda}}$. To see this, observe that 
$\tilde{\lambda}+\delta= ((n-1)(\ell+1)+\ell, (n-2)(\ell+1)+\ell, \ldots, (\ell+1)+\ell, \ell)$, 
which is also an arithmetic progression in which the difference between successive terms is $\ell+1$. Therefore, we have 
\vspace{0.4cm}
$$
a_{\tilde{\lambda}+\delta}=
\begin{vmatrix}x_1^\ell&x_2^\ell&\cdots&x_n^\ell\\x_1^{(\ell+1)+\ell}&x_2^{(\ell+1)+\ell}&\cdots&x_n^{(\ell+1)+\ell}\\\vdots&\vdots&\ddots&\vdots\\x_1^{(n-1)(\ell+1)+\ell}&x_2^{(n-1)(\ell+1)+\ell}&\cdots&x_n^{(n-1)(\ell+1)+\ell}\end{vmatrix}_{n\times n}
$$

\vspace{0.4cm}
$$
a_{\tilde{\lambda}+\delta}= \bigg{(} \prod_{i=1}^{n} x_i^{\ell}\bigg{)}
\begin{vmatrix}1&1&\cdots&1\\x_1^{\ell+1}&x_2^{\ell+1}&\cdots&x_n^{\ell+1}\\\vdots&\vdots&\ddots&\vdots\\x_1^{(n-1)(\ell+1)}&x_2^{(n-1)(\ell+1)}&\cdots&x_n^{(n-1)(\ell+1)}\end{vmatrix}_{n\times n}
$$  

\begin{align*}
s_{\tilde{\lambda}} &= \frac{\bigg{(} \prod_{i=1}^{n} x_i^{\ell}\bigg{)} \prod\limits_{\substack{i<j}} (x_j^{\ell+1}-x_i^{\ell+1})}{\prod\limits_{\substack{i<j}} (x_j-x_i)}\\
\end{align*}

In the above expression the numerator and denominator have small formulas and therefore using Lemma~\ref{lem:division of formulas} we get that $s_{\tilde{\lambda}}$ has a small formula of size $\poly(\ell,n)$.
\end{remark}

%
%
%
%
%

%

 


\subsubsection{Generalization to Skew Schur Polynomials}

A straightforward generalization of the previous result is to prove that a class of skew Schur polynomials is also hard for formulas. They can be defined via a Jacobi-Trudi like formula. 

\begin{theorem}\label{def: skew schur definition}
Let $\mu$ and $\lambda$ be partitions with $\mu_i \leq \lambda_i$ for every part $i$. Then, the skew Schur polynomial $s_{\lambda/\mu}$ satisfies

\begin{enumerate}
\item $s_{\lambda/\mu} = \det(h_{\lambda_i - \mu_j - i + j})_{1\leq i, j \leq k}$ where $k = l(\lambda).$
\item $s_{\lambda/\mu} = \det(e_{\lambda_i' - \mu_j' - i+j})_{1 \leq i, j \leq k}$ where $k = l(\lambda')$.  
\end{enumerate}
\end{theorem}

From these definitions of skew Schur polynomials, we can see that they also have ABPs of polynomial size, like Schur polynomials. However, skew Schur polynomials are in general linear combinations of Schur polynomials, by the Littlewood-Richardson rule $$s_{\lambda / \mu} = \sum_{\nu \vdash n - m} c^\lambda_{\mu, \nu} s_\nu.$$ The Littlewood-Richardson coefficients $c^\lambda_{\mu, \nu}$ count tableau whose Young diagram is of shape $\lambda / \mu$ and whose content satisfy certain technical conditions. They are also important in representation theory as they describe how Schur polynomials multiply, or equivalently how a tensor product of polynomial GL(n) representations decomposes into irreducible representations. They are also known to be \#P-hard to compute \cite{r29}. 

Hardness of skew Schur polynomials assuming hardness of determinant follows from the following lemma, Corollary~\ref{cor:prop s subsets} and Lemma~\ref{lem:key}.
%
%
%

%
%
%
%

%
\begin{lemma}\label{lem: skew distinct}
Let $l \geq 2$ and $\mu_1 \geq 1$ be positive integers. Let $\lambda, \mu$ be partitions with $l$ parts with $\lambda_i = (l-(i-1))l + \mu_1$ and $\mu_i = \mu_1$ for all $i < l$ and $\mu_l = \mu_1 - 1$. Then all entries of the homogeneous Jacobi-Trudi determinant $s_{\lambda/\mu}$ are distinct. 
\end{lemma}

\begin{proof}
For simplicity, we call the integer $k$ the label of the homogeneous symmetric polynomial $h_k$. Here, the $(i,j)$ entry of the matrix is $h_{(l-(i-1))l - i + j}$ for $1 \leq j < l$ and $h_{(l-(i-1))l + 1 - i + j}$ for $j = l$. Hence, for a fixed row $i$, all entries are distinct and the labels increase from left to right. Furthermore,

$$(l-i)l + 1 - (i+1) + l = (l-i)(l-1) + 1 - i < l(l-(i-1)) - i + 1$$

so the label of the last entry of row $(i+1)$ is strictly less than the label of the first entry of row $i$. Hence, all entries of the Jacobi-Trudi determinant are distinct.
\end{proof}

%
%

\begin{thebibliography}{41}
\bibitem{r1} Markus Bläser and Gorav Jindal. On the Complexity of Symmetric Polynomials. In Avrim Blum, editor, 10th Innovations in Theoretical Computer Science Conference (ITCS 2019), volume 124 of Leibniz International Proceedings in Informatics (LIPIcs), pages 47:1–47:14, Dagstuhl, Germany, 2018. Schloss Dagstuhl–Leibniz-Zentrum fuer Informatik. doi:10.4230/ LIPIcs.ITCS.2019.47.
\bibitem{r2} Peter Bürgisser. Completeness and Reduction in Algebraic Complexity Theory. Algorithms and Computation in Mathematics. Springer, 2000. URL: https://www.springer.com/gp/book/ 9783540667520.
\bibitem{r3} Cy Chan, Vesselin Drensky, Alan Edelman, Raymond Kan, and Plamen Koev. On computing schur functions and series thereof. Journal of Algebraic Combinatorics, 50(2):127–141, September 2019. doi:10.1007/s10801-018-0846-y.
\bibitem{r4} Xi Chen, Neeraj Kayal, and Avi Wigderson. Partial derivatives in arithmetic complexity. Foundations and Trends in Theoretical Computer Science, 2011.
\bibitem{r5} Chi-Ning Chou, Mrinal Kumar, and Noam Solomon. Some closure results for polynomial factorization and applications. CoRR, abs/1803.05933, 2018. arXiv:1803.05933.
\bibitem{r6} Chi-Ning Chou, Mrinal Kumar, and Noam Solomon. Closure of VP under taking factors: a short and simple proof. CoRR, abs/1903.02366, 2019. arXiv:1903.02366.
\bibitem{r7} James Demmel and Plamen Koev. Accurate and efficient evaluation of schur and jack functions. Mathematics of Computation, 75(253):223–239, 2006. URL: http://www.jstor.org/stable/ 4100151.
\bibitem{r8} Pranjal Dutta, Nitin Saxena, and Amit Sinhababu. Discovering the roots: Uniform closure results for algebraic classes under factoring. In STOC, pages 1152–1165, 2018. doi:10.1145/ 3188745.3188760.
\bibitem{r9} Zeev Dvir, Amir Shpilka, and Amir Yehudayoff. Hardness-randomness tradeoffs for bounded depth arithmetic circuits. SIAM J. Comput., 39(4):1279–1293, 2010. Preliminary version in STOC’08. doi:10.1137/080735850.
\bibitem{r10} Philippe Flajolet and Robert Sedgewick. Analytic Combinatorics. Cambridge University Press, 2014.
\bibitem{r11} Sergey Fomin, Dima Grigoriev, and Gleb Koshevoy. Subtraction-free complexity, cluster transformations, and spanning trees. Found. Comput. Math., 16(1):1–31, February 2016. doi:10.1007/s10208-014-9231-y.
\bibitem{r12} Sergey Fomin, Dima Grigoriev, Dorian Nogneng, and Éric Schost. On semiring complexity of schur polynomials. Comput. Complex., 27(4):595–616, December 2018. doi:10.1007/ s00037-018-0169-3.
\bibitem{r13} Hervé Fournier, Nutan Limaye, Meena Mahajan, and Srikanth Srinivasan. The shifted partial derivative complexity of elementary symmetric polynomials. Theory of Computing, 13(1):1–34, 2017. doi:10.4086/toc.2017.v013a009.
\bibitem{r14} Vesselin Gasharov. Incomparability graphs of (3+ 1)-free posets are s-positive. Discrete Mathematics, 157(1-3):193–197, 1996.
\bibitem{r15} Dima Grigoriev and Gleb Koshevoy. Complexity of tropical schur polynomials. Journal of Symbolic Computation, 74:46–54, 2016. doi:10.1016/j.jsc.2015.05.005.
\bibitem{r16} Ankit Gupta, Pritish Kamath, Neeraj Kayal, and Ramprasad Saptharishi. Arithmetic circuits: A chasm at depth 3. SIAM J. Comput., 45(3):1064–1079, 2016. doi:10.1137/140957123.
\bibitem{r17} Sean Hallgren, Alexander Russell, and Amnon Ta-Shma. Normal subgroup reconstruction and quantum computation using group representations. In Proceedings of the thirty-second annual ACM symposium on Theory of computing, pages 627–635. ACM, 2000.
\bibitem{r18} E. R. Heineman. Generalized vandermonde determinants. Transactions of the American Mathematical Society, 31(3):464–476, 1929. URL: http://www.jstor.org/stable/1989528.
\bibitem{r19} Pavel Hrubes and Amir Yehudayoff. Homogeneous formulas and symmetric polynomials. Computational Complexity, 20(3):559–578, 2011. doi:10.1007/s00037-011-0007-3.
\bibitem{r20} Christian Ikenmeyer and Stefan Mengel. On the relative power of reduction notions in arithmetic circuit complexity. Information Processing Letters, 130:7–10, 2018.
\bibitem{r21} Christian Ikenmeyer and Greta Panova. Rectangular kronecker coefficients and plethysms in geometric complexity theory. Advances in Mathematics, 319:40–66, 2017.
\bibitem{r22} Valentine Kabanets and Russell Impagliazzo. Derandomizing polynomial identity tests means proving circuit lower bounds. Computational Complexity, 13(1-2):1–46, 2004. Preliminary version in the 35th Annual ACM Symposium on Theory of Computing (STOC 2003). doi: 10.1007/s00037-004-0182-6.
\bibitem{r23} Erich Kaltofen. Factorization of Polynomials Given by Straight-Line Programs. In Randomness and Computation, pages 375–412. JAI Press, 1989.
\bibitem{r24} Plamen. Koev. Accurate computations with totally nonnegative matrices. SIAM Journal on Matrix Analysis and Applications, 29(3):731–751, 2007. doi:10.1137/04061903X.
\bibitem{r25} J. M. Landsberg and Zach Teitler. On the ranks and border ranks of symmetric tensors. Foundations of Computational Mathematics, 10(3):339–366, June 2010. doi:10.1007/ s10208-009-9055-3.
\bibitem{r26} Veerle Ledoux and Simon JA Malham. Introductory schubert calculus, 2010.
\bibitem{r27} Dick Lipton and Ken Regan. Arithmetic complexity and symmetry, 2009. URL: https: //rjlipton.wordpress.com/2009/07/10/arithmetic-complexity-and-symmetry/.
\bibitem{r28} Ian G. Macdonald. Symmetric functions and Hall polynomials. Oxford University Press, 1979.
\bibitem{r29} Hariharan Narayanan. On the complexity of computing kostka numbers and littlewoodrichardson coefficients. Journal of Algebraic Combinatorics, 24(3):347–354, 2006.
\bibitem{r30} Noam Nisan and Avi Wigderson. Lower bounds on arithmetic circuits via partial derivatives. Computational Complexity, 6(3):217–234, 1997. Available on citeseer:10.1.1.90.2644. doi: 10.1007/BF01294256.
\bibitem{r31} Ryan O’Donnell and John Wright. Quantum spectrum testing, 2015. arXiv:1501.05028.
\bibitem{r32} Luke Oeding. Border ranks of monomials, 2016. arXiv:1608.02530.
\bibitem{r33} Bruce Sagan. The Symmetric Group: Representations, Combinatorial Algorithms, and Symmetric Functions. Springer, 2001.
\bibitem{r34} Nitin Saxena. Diagonal Circuit Identity Testing and Lower Bounds. In Proceedings of the 35th International Colloquium on Automata, Languages and Programming (ICALP 2008), pages 60–71, 2008. doi:10.1007/978-3-540-70575-8\_6.
\bibitem{r35} Amir Shpilka. Affine projections of symmetric polynomials. In In Proc. 16th Annual IEEE Conference on Computational Complexity, pages 160–171, 2001.
\bibitem{r36} Amir Shpilka and Avi Wigderson. Depth-3 arithmetic circuits over fields of characteristic zero. Computational Complexity, 10(1):1–27, 2001. Preliminary version in the 14th Annual IEEE Conference on Computational Complexity (CCC 1999). doi:10.1007/PL00001609.
\bibitem{r37} Amir Shpilka and Amir Yehudayoff. Arithmetic circuits: A survey of recent results and open questions. Foundations and Trends in Theoretical Computer Science, 5:207–388, March 2010. doi:10.1561/0400000039.
\bibitem{r38} Richard P Stanley. On the number of reduced decompositions of elements of coxeter groups. European Journal of Combinatorics, 5(4):359–372, 1984.
\bibitem{r39} Richard P Stanley. Enumerative combinatorics. vol. 2, volume 62 of. Cambridge Studies in Advanced Mathematics, 1999.
\bibitem{r40} Leslie G. Valiant, Sven Skyum, S. Berkowitz, and Charles Rackoff. Fast Parallel Computation of Polynomials Using Few Processors. SIAM J. Comput., 12(4):641–644, 1983. Preliminary version in the 6th Internationl Symposium on the Mathematical Foundations of Computer Science (MFCS 1981). doi:10.1137/0212043.
\bibitem{r41} Herbert S. Wilf. Generatingfunctionology. A. K. Peters, Ltd., third edition, 2006.
\end{thebibliography}
\end{document}
